%%%%%%%%%%%%%%%%%%%%%%%%%%%%%%%%%%%%%%%%%%%%%%%%%%%%%%%%%%%%%%%%%%%%
%% I, the copyright holder of this work, release this work into the
%% public domain. This applies worldwide. In some countries this may
%% not be legally possible; if so: I grant anyone the right to use
%% this work for any purpose, without any conditions, unless such
%% conditions are required by law.
%%%%%%%%%%%%%%%%%%%%%%%%%%%%%%%%%%%%%%%%%%%%%%%%%%%%%%%%%%%%%%%%%%%%

\documentclass[
  digital,     %% The `digital` option enables the default options for the
               %% digital version of a document. Replace with `printed`
               %% to enable the default options for the printed version
               %% of a document.
%%  color,       %% Uncomment these lines (by removing the %% at the
%%               %% beginning) to use color in the printed version of your
%%               %% document
  oneside,     %% The `oneside` option enables one-sided typesetting,
               %% which is preferred if you are only going to submit a
               %% digital version of your thesis. Replace with `twoside`
               %% for double-sided typesetting if you are planning to
               %% also print your thesis. For double-sided typesetting,
               %% use at least 120 g/m² paper to prevent show-through.
  nosansbold,  %% The `nosansbold` option prevents the use of the
               %% sans-serif type face for bold text. Replace with
               %% `sansbold` to use sans-serif type face for bold text.
  nocolorbold, %% The `nocolorbold` option disables the usage of the
               %% blue color for bold text, instead using black. Replace
               %% with `colorbold` to use blue for bold text.
  lof,         %% The `lof` option prints the List of Figures. Replace
               %% with `nolof` to hide the List of Figures.
  lot,         %% The `lot` option prints the List of Tables. Replace
               %% with `nolot` to hide the List of Tables.
]{fithesis4}
%% The following section sets up the locales used in the thesis.
\usepackage[resetfonts]{cmap} %% We need to load the T2A font encoding
\usepackage[T1,T2A]{fontenc}  %% to use the Cyrillic fonts with Russian texts.
\usepackage[
  main=english, %% By using `czech` or `slovak` as the main locale
                %% instead of `english`, you can typeset the thesis
                %% in either Czech or Slovak, respectively.
  english, german, czech, slovak %% The additional keys allow
]{babel}        %% foreign texts to be typeset as follows:
%%
%%   \begin{otherlanguage}{german}  ... \end{otherlanguage}
%%   \begin{otherlanguage}{czech}   ... \end{otherlanguage}
%%   \begin{otherlanguage}{slovak}  ... \end{otherlanguage}
%%
%%
%% The following section sets up the metadata of the thesis.
\thesissetup{
    date        = \the\year/\the\month/\the\day,
    university  = mu,
    faculty     = fi,
    type        = mgr,
    department  = Department of Machine Learning and Data Processing,
    author      = Jakub Čieško,
    gender      = m,
    advisor     = {doc. RNDr. Aleš Horák, Ph.D.},
    title       = {Scene-Aware Dialogue for the Pepper Robot Using Object Tracking and Large Language Models},
    TeXtitle    = {Scene-Aware Dialogue for the Pepper Robot Using Object Tracking and Large Language Models},
keywords    = {Social Robotics, Large Language Models, Vision-Language Models, Scene Graph Generation, Object Tracking, Object Detection, RAG, Pepper Robot},
    TeXkeywords = {Social Robotics, Large Language Models, Vision-Language Models, Scene Graph Generation, Object Tracking, Object Detection, RAG, Pepper Robot},
    abstract    = {%
      The goal of the thesis is to design and implement a dialogue interface system for natural communication with the Pepper robot, based on its visual perception of its surroundings. The student will build upon an existing object detection solution and extend it with multi-frame video object tracking and the construction of a dynamic semantic scene model. The work will include detecting spatial relationships between objects (such as occlusion or relative position) and integrating these findings with the robot's internal modules (specifically ALPeoplePerception). The resulting structured representation of the environment will serve as a dynamic information source for the dialogue manager. The dialogue system will be implemented using Large Language Models (LLMs) and Retrieval-Augmented Generation (RAG) techniques, where the model uses the current scene context to conduct fluid and factually supported conversation about the objects and persons the robot sees in its surroundings.

The final submission will consist of a textual part with overviews of related technologies and approaches and a detailed description of the implementation, evaluation and testing, as well as a working implementation of vision recognition and dialog with the Pepper robot about its surroundings.

    },
    thanks      = {%
      I would not have been able to complete this thesis without the guidance of Doc. Horák, whose supervision provided both valuable direction and the freedom to explore methods of my own interest. I would also like to thank Dr. Rambousek for his patience during the robot testing process. Finally, I am deeply grateful to my close ones for their support and understanding throughout this particularly demanding period.
    },
    bib         = bibliography_clean.bib,
    %% Remove the following line to use the JVS 2018 faculty logo.
    facultyLogo = fithesis-fi,
}
\usepackage{makeidx}      %% The `makeidx` package contains
\makeindex                %% helper commands for index typesetting.
%% These additional packages are used within the document:
\usepackage{paralist} %% Compact list environments
\usepackage{amsmath}  %% Mathematics
\usepackage{amsthm}
\usepackage{amsfonts}
\usepackage{url}      %% Hyperlinks
\usepackage{markdown} %% Lightweight markup
\usepackage{listings} %% Source code highlighting
\lstset{
  basicstyle      = \ttfamily,
  identifierstyle = \color{black},
  keywordstyle    = \color{blue},
  keywordstyle    = {[2]\color{cyan}},
  keywordstyle    = {[3]\color{olive}},
  stringstyle     = \color{teal},
  commentstyle    = \itshape\color{magenta},
  breaklines      = true,
}
\usepackage{floatrow} %% Putting captions above tables
\floatsetup[table]{capposition=top}
\usepackage[babel]{csquotes} %% Context-sensitive quotation marks



\begin{document}
%% The \chapter* command can be used to produce unnumbered chapters:
%%\chapter*{Introduction}
%% Unlike \chapter, \chapter* does not update the headings and does not
%% enter the chapter to the table of contents. I we want correct
%% headings and a table of contents entry, we must add them manually:
%\markright{\textsc{Introduction}}
%\addcontentsline{toc}{chapter}{Introduction}

\chapter{Introduction}

The motivation behind this thesis as well as the research problem and overall direction is introduced in this chapter. 
First, the work is situated within the broader field of human-robot interaction (HRI) and the growing role of social robots in real-world environments. 
It then identifies the central limitation of current conversational robotic systems, namely the lack of grounding in perceived context, and motivates the need for integrating perception, memory, and language into a unified scene-aware dialogue framework. 
Based on this motivation, the research objectives and questions are formulated at the end of the chapter.

\section{Motivation}

Social robots, that is, robots designed for interaction, communication, and assistance in human-centered environments, are becoming increasingly present in everyday life. They are already being deployed in domains such as healthcare, assisted living, education, public services, and industry. As their presence expands, social robotics is no longer only an experimental research topic for the labs, but an increasingly practical technological reality. This broader adoption also raises a fundamental question: how can robots operate meaningfully in environments shaped by human needs, expectations, and social norms?

This question places human-robot interaction (HRI) at the center of contemporary robotics research. Sheridan's well-known taxonomy distinguishes supervisory control, teleoperation, automated vehicles, and social interaction~\cite[525--526]{Sheridan2016HRI}. This thesis focuses on the latter category, in which robots are expected to assist, guide, teach, comfort, or accompany users through direct engagement. While the emphasis here is on social interaction, the underlying challenge extends beyond a single HRI category. Once robots leave tightly controlled industrial settings and enter shared human environments, successful operation increasingly depends not only on mechanical capability, but also on the quality of interaction they can sustain with users.

The importance of this challenge is especially visible in healthcare and education, where interaction quality is often inseparable from task success. In healthcare, HRI is inherently interdisciplinary, combining engineering, social sciences, and design to create systems that can support people in meaningful and acceptable ways. Robots are being explored as monitoring tools, rehabilitation assistants, guides, and providers of emotional support~\cite{Han2024HRIHealthcare}. Applications further extend to such domains as autism therapy and educational settings, where robots such as NAO and Pepper by Aldebaran, two platforms widely used in HRI research, have shown promising, although still mixed, results~\cite{Tapus2012AutismNao, aturhurra_leveraging_2024}. Pepper has also been studied as an intelligent assistant in libraries using visual perception and OCR, while multimodal assistants combining vision and large language models have been proposed for environments such as physics laboratories~\cite{trinquet_future_2025, latif_physicsassistant_2024}. Across these examples, a common conclusion emerges: the usefulness of a social robot does not arise from speech alone or from any isolated subsystem, but from its ability to connect interaction to its surrounding situation.

This requirement naturally leads to the concept of grounding. For interaction to be meaningful, a robot should not merely generate fluent responses, but relate them to what it currently perceives. In practice, this includes identifying relevant objects and people, recognizing attributes and relationships, and responding in a way that is consistent with the observed environment. Without such grounding, conversational ability can create only an illusion of understanding rather than genuine situational competence.

Recent advances in large language models (LLMs) have made this issue even more important. LLM-based systems can replace traditionally modular pipelines composed of automatic speech recognition, intent detection, dialogue management, and text-to-speech with a more unified generative framework~\cite{aturhurra_leveraging_2024}. This enables more flexible and natural interaction, but it also increases user expectations that the robot genuinely understands the current situation. Even though LLM-based systems can substantially improve conversational richness and engagement, fluent language without perceptual grounding creates a mismatch between apparent competence and real-world understanding~\cite{kim_understanding_2024}. In socially assistive settings, this mismatch also raises broader concerns related to trust, bias and requires more responsible deployment~\cite{aturhurra_leveraging_2024}.

Achieving grounded interaction, however, requires more than attaching a language model to camera input. Perception and dialogue operate at different levels of abstraction and therefore need an intermediate representational layer. A structured scene representation, such as a scene graph describing objects, their attributes, and their relationships, provides such a bridge between low-level visual outputs and higher-level reasoning~\cite{johnson_image_2015}. Structured representations support relational queries, more precise references, and more coherent communication about the environment.

A further requirement is persistence over time. Real-world interaction is rarely limited to a single observation; it unfolds across multiple conversation turns and changing scenes. A robot should thus maintain a form of memory that preserves object identity across observations, stores relevant attributes and relationships, and supports questions not only about the current scene but also about past context. This aligns with broader trends in robotics, where higher-level competence is increasingly associated with persistent and semantically rich scene representations rather than purely geometric or single frame based perception~\cite{mascaro_scene_2025, peng_object_2024}.

For the Pepper robot, these requirements are accompanied by an additional practical constraint: limited onboard computational resources. State-of-the-art models for detection, captioning, multimodal reasoning, and dialogue generation are computationally demanding and cannot be executed efficiently on Pepper alone\footnote{One should note that running a lightweight detection could in principle be possible even on not so computationally powerful hardware; but a full pipeline composed of detection, language modeling, and other steps would be impossible to run on edge devices of this capacity.}. Therefore, many modern robotic systems rely on distributed architectures in which perception and reasoning are offloaded to external compute resources. Prior work has demonstrated near real-time perception through external edge hardware~\cite{reyes_near_2018, wilson_enhancing_2025}, distributed industrial assistants combining local and remote resources~\cite{becerra_improving_2024, brad_retrieval-augmented_2025, hafez_enhancing_nodate}, and visually grounded conversational systems that combine Pepper with external captioning and language models~\cite{grassi_grounding_2024, latif_physicsassistant_2024, trinquet_future_2025}. Off-board computation is therefore not merely an implementation convenience, but a central architectural enabler for advanced capabilities on computationally constrained robot platforms.

The motivation of this thesis is then to move beyond isolated perception or language components toward a unified system for scene-aware dialogue. In this thesis, scene-aware dialogue denotes the ability of a robot to perceive its environment, identify and track relevant entities, represent their attributes and relationships, maintain this information over time, and use it as a basis for coherent grounded conversation. By integrating modern computer vision, structured scene representations, persistent memory, and conversational models within a distributed client-server architecture, the thesis aims to advance social robots toward more context-aware, reliable, and meaningful interaction in real-world environments.

\section{Problem Statement and Objectives}
The motivation outlined in the previous section identified a central gap in contemporary human-robot interaction: conversational ability has advanced rapidly, yet robust interaction still requires explicit connection to perception, memory, and situational context. This thesis addresses that gap through the design of a unified scene-aware interaction framework for the Pepper robot.

More specifically, the central problem is how to enable Pepper to perceive its surroundings, transform perceptual observations into a structured and persistent internal representation, and use this evolving world model as the basis for coherent dialogue with the user. In practical terms, this requires the robot to maintain awareness of relevant objects and people, preserve identity across observations, incorporate robot-native sensing signals, and support context-sensitive follow-up questions grounded in both the current scene and recent interaction history.

Addressing this problem is then not only a matter of dialogue generation, but it also requires solving a broader systems challenge of connecting heterogeneous components, namely the robot sensors, external (multimodal) models, scene memory, and language generation, into a responsive and reliable architecture suitable for real-world interaction.

To address these requirements, the thesis defines the following objectives:

\begin{itemize}
    \item Design a distributed system architecture that enables the use of modern perception and language models despite the computational limitations of the robot platform.
    
    \item Develop a perception pipeline that preserves object identity over time.
    
    \item Construct a structured scene representation that captures entities, their attributes, and their relationships.
    
    \item Integrate server-side perception with robot-native sensing (e.g., NAOqi APIs), particularly for human-related and social signals.
    
    \item Enable grounded dialogue based on current and recent scene state, including support for multi-turn interaction and multilingual communication in Czech and English.
    
    \item Provide tools for monitoring, configuration, and evaluation, including assessment of scene representation quality and overall system behavior.
\end{itemize}

The proposed solution does not however exist in isolation. It builds on several strands of prior work. 
The general idea of grounding dialogue in visual scene understanding is inspired by the system introduced by Grassi et al.~\cite{grassi_grounding_2024}. The use of external computation to overcome Pepper's hardware limitations follows established approaches based on edge or remote processing~\cite{reyes_near_2018, trinquet_future_2025, brad_retrieval-augmented_2025}. Additional design choices draw on a broader body of related research, which is discussed in detail in the corresponding chapters of this thesis.

\section{Research Questions (MAYBE REDUNDANT SECTION, ALREDY PRESENT IN OBJECTIVES?)}
(MAYBE REDUNDANT SECTION, ALREDY PRESENT IN OBJECTIVES? I just wanted research questions in the thesis.)
The problem formulation leads to these five research questions. 
\begin{enumerate}
    \item \textbf{System architecture:} How can Pepper's limited onboard capabilities be extended through a remote multimodal inference server while preserving responsiveness and maintaining grounded interaction?

    \item \textbf{Scene representation:} How should object detection, persistent tracking, and scene graph generation be combined to form a dynamic world model suitable for dialogue?

    \item \textbf{Multimodal fusion:} What benefits, if any, arise from integrating Pepper-native people perception and social cues with server-side visual detections?

    \item \textbf{Scene graph methods:} To what extent can a hybrid scene graph approach, combining rule-based reasoning, vision-language model prompting, and learned relation prediction, support grounded dialogue?

    \item \textbf{Robustness:} How sensitive is scene graph quality to factors such as prompt design, vocabulary size, and model or provider choice?
\end{enumerate}

\section{Main Contributions}
This thesis tries to make several contributions at the intersection of social robotics, multimodal AI, and grounded human-robot interaction.

The primary contribution of this thesis is the design and implementation of a scene-aware dialogue system for the Pepper robot that integrates embodied interaction with server-side multimodal inference and explicit scene memory.

A second contribution is the development of a configurable server platform based on FastAPI that supports multimodal inference, runtime configuration, and persistent scene state management. The platform enables multiple interaction modes, including captioning, object detection, scene graph generation, and memory-grounded dialogue, while remaining adaptable to different models, providers, and deployment settings. It further incorporates an optional RPC-based worker architecture that isolates GPU workloads and enables efficient allocation, restart, and release of VRAM resources.

A further contribution is the design of a unified perception pipeline that combines object detection, persistent tracking, Set-of-Marks prompting for hybrid scene graph generation, captioning, and memory updates within a single processing flow. As part of this pipeline, a Pepper-specific fusion layer aligns server-side detections with robot-native sensing, particularly for human-related and social signals.

Moreover, this thesis also contributes a dialogue layer that leverages structured scene memory and recent perceptual context to support grounded interaction through retrieval-augmented generation (RAG). This layer enables general scene-aware dialogue, object-focused queries, attribute-focused queries, relationship-focused queries, cached pregenerated question-answer pairs, and multilingual interaction in Czech and English.

Another contribution is the supporting infrastructure for inspection, operation, and evaluation of the system. This includes a browser-based dashboard for live monitoring, configuration, memory inspection, and system control, together with an experimental framework for scene description generation, vocabulary construction, scene graph analysis, and comparative evaluation.

On the robot side, the thesis develops a modular Pepper client responsible for image acquisition, robot-context collection using NAOqi APIs, interaction orchestration, speech output, tablet interaction, and communication with the external server. The client bridges Pepper's native capabilities, such as QiChat grammar, with modern external AI services while preserving an interactive user experience.

Taken together, these contributions establish both a functional Pepper-based prototype and a reusable platform for studying grounded dialogue in environments.

\section{Thesis Structure}

The remainder of the thesis is organized as follows.

Chapter~2 reviews the relevant theoretical background and related work in social robotics, grounded dialogue, visual perception, scene representations, and multimodal AI systems. It establishes the conceptual and technical foundations on which the proposed system is built.

Chapter~3 presents the overall system architecture and design of the proposed framework. It introduces the main components of the system, their responsibilities, the interaction flow between them, and the key design decisions arising from the requirements identified in the introductory chapter.

Chapter~4 focuses on the server-side implementation. It describes the multimodal inference platform, perception pipeline, scene graph generation methods, scene memory, dialogue services, configuration mechanisms, and supporting runtime infrastructure.

Chapter~5 describes the robot-side integration on the Pepper platform. It covers image acquisition, robot-context collection, interaction orchestration, speech and dialogue integration, tablet interface components, and communication with the external server.

Chapter~6 presents the experimental methodology and evaluation. It focuses on Set-of-Marks prompting strategies, scene graph generation approaches, prompt design, and related factors influencing structured scene understanding quality as well as system latency.

Chapter~7 brings together the final evaluation of the proposed system through results, discussion, and concluding remarks. It summarizes the main findings, discusses limitations and trade-offs, and outlines directions for future work.

The appendices provide supplementary technical material, including prompts, structured output schemas, API definitions, configuration examples, selected interaction transcripts, and additional experimental details.

\chapter{Background and Related Work}

This chapter reviews the theoretical foundations and related work necessary for building a scene-aware dialogue system. It progresses from the high-level requirements of human-robot interaction down to the specific visual perception models, structured scene representations, and efficient multimodal language generation techniques utilized in this thesis.

\section{Social Robotics and Grounded Human-Robot Interaction (HRI)}

Social robots operate in shared human environments where users expect them to exhibit awareness of their physical and social surroundings. Recent integrations of Large Language Models (LLMs) into human-robot interaction (HRI) have significantly improved conversational fluency. However, as Kim et al.~\cite{kim_understanding_2024} and Atuhurra et al.~\cite{atuhurra_leveraging_2024} note, this fluency often raises user expectations: when a robot speaks naturally, users implicitly assume it also understands the physical context and social nuances of the conversation.

To prevent a mismatch between linguistic competence and actual situational awareness, the robot's dialogue must be grounded. In robotics, grounding denotes the explicit connection between linguistic expressions (e.g., ``the person on my left'') and the physical entities to which they refer. Unlike a purely text-based assistant, an embodied robot must continuously resolve these references against its changing perceptual state.

The Pepper robot is a prominent platform for HRI research due to its humanoid form, native tablet interface~\cite{softbank_peppers_tablet_2026}, and built-in speech recognition~\cite{softbank_alspeechrecognition_2026}. However, its onboard computational hardware is severely limited and insufficient for modern deep learning inference. Consequently, achieving near real-time scene understanding on Pepper requires a distributed architecture, where heavy visual and linguistic processing is offloaded to an external server while the robot handles interaction and sensory acquisition~\cite{reyes_near_2018, trinquet_future_2025, brad_retrieval-augmented_2025}.

\section{Visual Perception for Robotics}

The foundation of scene awareness is object detection. Traditionally, robotics has relied on Convolutional Neural Networks (CNNs), such as the YOLO family, which utilize anchor boxes and dense grid predictions. While fast, these models require heuristic post-processing like Non-Maximum Suppression (NMS). Recently, the architectural shift toward Detection Transformers (DETRs)~\cite{zhao_detrs_2024} has reframed detection as a direct set-prediction problem. By leveraging self-attention mechanisms~\cite{vaswani2023attentionneed} and bipartite matching during training, models like RT-DETR and RF-DETR eliminate NMS, providing state-of-the-art accuracy while maintaining the real-time latency required for interactive robotics.

Standard detectors operate on a closed vocabulary, mapping visual features to a fixed set of classes via a classifier head. This restricts natural dialogue, as the robot cannot recognize novel objects. Open-vocabulary models, such as OWL-ViT~\cite{minderer_simple_2022} and Grounding DINO~\cite{liu_grounding_2024}, break this barrier by projecting image patches and text prompts into a shared multimodal embedding space (often aligned via CLIP~\cite{pmlr-v139-radford21a}). This allows the detection of arbitrary objects based on text similarity, a crucial capability for open-ended scene-aware dialogue.

Detecting an object in a single frame is insufficient for conversation; the system must preserve object permanence. Multi-Object Tracking (MOT) algorithms associate detections across frames using geometric overlap (IoU) and appearance embeddings~\cite{bewley_simple_2016, wojke_simple_2017, zhang_bytetrack_2022}. In a dialogue context, robust tracking allows the robot to reliably answer follow-up questions about an entity, sustaining referential continuity.

\section{Semantic Scene Understanding and Scene Graphs}

While bounding boxes localize objects, they lack the relational semantics required for language generation. Scene Graphs~\cite{johnson_image_2015} bridge this gap by providing a structured intermediate representation---a directed graph where nodes represent detected objects and edges represent semantic or spatial relationships (e.g., \verb|<cup> [on] <table_>|).

In addition to pairwise spatial relations, objects possess unary properties such as color or state. In this work, to maintain a uniform graph structure, attributes are elegantly modeled as self-referential edges (unary relations) on the object nodes. This unifies the representation, allowing the dialogue system to query properties and spatial relations using identical graph traversal logic.

Scene Graph Generation (SGG) can be achieved through heuristic geometric rules or learned models. Learned approaches, such as RelTR (Relation Transformer)~\cite{cong_reltr_2023}, treat SGG as a direct set-prediction problem, inferring subject-predicate-object triplets from image features. Reflecting the consensus that hierarchical, persistent representations are central to robust robot understanding~\cite{mascaro_scene_2025, hughes_foundations_2024}, this thesis employs a hybrid neuro-symbolic approach, combining fast geometric rules for stable spatial relations with learned multimodal prompting for richer semantic connections.

\section{Multimodal AI: LLMs, VLMs, and Efficient Architectures}

To generate grounded dialogue, the structured scene representation must be connected to natural language via Large Language Models (LLMs). LLMs are built upon the Transformer architecture, utilizing self-attention mechanisms for autoregressive next-token prediction~\cite{vaswani2023attentionneed}. 

Vision-Language Models (VLMs) extend this capability by conditioning text generation on visual inputs~\cite{liu_visual_2023, dai_instructblip_2023, driess_palm-e_2023}. Typically, a Vision Encoder (such as a ViT) extracts features from an image, which are then projected into the LLM's embedding space. To the LLM, the image functions as a sequence of ``visual words'' appended to the text prompt.

As models scale to achieve better reasoning, computational costs rise dramatically. A dense model activates all its parameters for every token. To mitigate this, modern architectures frequently employ a Mixture of Experts (MoE) routing mechanism~\cite{shazeer2017outrageouslylargeneuralnetworks}. In an MoE model, only a small subset of specialized neural ``experts'' is activated for any given token. Consequently, if a system possesses sufficient RAM or VRAM to load the model weights, it can achieve the reasoning capabilities of a massive parameter model at the inference speed and compute cost of a much smaller model. This architectural efficiency is highly advantageous in robotics, where rapid inference is critical.

Despite their capabilities, VLMs are prone to hallucinating relations or failing to disambiguate multiple instances of the same object. To enforce precise visual grounding, this thesis utilizes Set-of-Mark (SoM) prompting~\cite{yang_set-of-mark_2023}. By rendering explicit visual markers---such as numbered bounding boxes or segmentation masks~\cite{kirillov_segment_2023}---directly onto the image, SoM transforms visual regions into explicit, referable symbols. This forces the VLM to anchor its relational reasoning to specific, tracked detections.

\section{Grounded Dialogue and Retrieval-Augmented Generation (RAG)}

Classical robotic dialogue systems rely on rigid intent-slot architectures~\cite{bocklisch_rasa_2017}. While predictable, these rule-based managers become brittle when users ask unanticipated questions about the environment. Generative AI offers open-ended flexibility, provided the generation can be tethered to the robot's physical reality.

Retrieval-Augmented Generation (RAG)~\cite{lewis_retrieval-augmented_2020} provides a framework for enriching language models with external data. While traditional RAG retrieves static text passages from document databases, this thesis adapts the paradigm for embodied robotics. Here, the ``database'' is the robot's dynamic Scene Memory. The dialogue layer retrieves structured scene graphs, recent visual captions, and object-specific attributes, injecting this live perceptual context directly into the LLM's prompt~\cite{grassi_grounding_2024}.

Crucially, this structured RAG approach shifts the cognitive burden placed on the LLM. Rather than relying on its internal Parametric Memory to hallucinate an answer, the model is tasked with ``reading comprehension'' over the injected prompt context. By providing the exact physical state in the prompt, even relatively small LLMs can perform exceptionally well on grounded Question-Answering tasks. Coupled with the MoE architectures discussed previously, this strategy drastically reduces VRAM requirements and inference latency, enabling highly responsive, factually grounded HRI on computationally constrained platforms.


\chapter{System Architecture and Design}

\section{Use Cases and Interaction Scenarios}
\begin{itemize}
    \item Define the concrete user-facing interaction modes implemented by the system: \texttt{look}, \texttt{scan}, \texttt{ask}, \texttt{refreshAndAsk}, and conversation reset.
    \item Describe typical scenarios such as single-frame captioning, multi-frame scan summaries, object-oriented follow-up questions, and direct vision chat over an uploaded image.
    \item State which scenarios are evaluated offline and which are evaluated in embodied robot interaction.
\end{itemize}

\section{Functional Requirements}
\begin{itemize}
    \item Remote image processing from Pepper.
    \item Configurable object detection backends.
    \item Persistent object identities across frames.
    \item Scene graph generation with attributes and relations.
    \item Integration of Pepper-native people and social perception into the world model.
    \item Grounded dialogue based on current scene memory and recent conversation history.
    \item Multilingual output support.
    \item Runtime configuration, monitoring, and memory inspection through a dashboard.
    \item Offline experimental tooling for prompt, vocabulary, and graph-quality studies.
\end{itemize}

\section{Non-Functional Requirements and Constraints}
\begin{itemize}
    \item Pepper-side hardware and software limitations.
    \item Network dependence and the need for graceful degradation.
    \item Latency constraints for interactive dialogue.
    \item Modularity and provider interchangeability.
    \item Traceability of intermediate states for debugging and evaluation.
    \item Safety and robustness under partial sensor failure, provider failure, or malformed structured outputs.
\end{itemize}

\section{Design Principles}
\begin{itemize}
    \item Separate robot-side sensing and interaction from server-side heavy inference.
    \item Keep state explicit and inspectable through a structured scene memory.
    \item Prefer configurable components over hard-coded pipelines.
    \item Support both fast approximate reasoning and slower richer reasoning.
    \item Reuse production components inside the research pipeline where possible.
\end{itemize}

\section{System Boundaries}
\begin{itemize}
    \item Clarify what is in scope: perception, memory, grounded dialogue, runtime configuration, and evaluation tooling.
    \item Clarify what is not fully implemented or not central: full autonomous action planning, navigation, long-term vector memory, and broad manipulation-oriented robot control.
\end{itemize}


\section{End-to-End Processing Pipeline}
\begin{itemize}
    \item Present the complete data flow from Pepper image capture to spoken answer.
    \item Include the main stages: capture, metadata collection, HTTP request, server-side inference, memory update, dialogue generation, and robot speech output.
    \item Add a figure with the full end-to-end architecture.
\end{itemize}

\section{Deployment Topology}
\begin{itemize}
    \item Describe the three main runtime domains: Pepper client, FastAPI orchestration server, and optional worker process.
    \item Explain when inference runs in-process and when it is delegated to the worker runtime.
    \item Discuss external providers, local GPU models, and optional tunnel exposure for remote access.
\end{itemize}

\section{Communication Interfaces and Data Contracts}
\begin{itemize}
    \item Describe the REST endpoints for captioning, detection, chat, vision chat, configuration, worker control, and memory access.
    \item Explain the robot metadata payload and why it is essential for person fusion and geometric reasoning.
    \item Describe the WebSocket channel used by the dashboard for live updates.
\end{itemize}

\section{Runtime Configuration and Reconfiguration}
\begin{itemize}
    \item Describe the YAML plus Pydantic configuration model.
    \item Explain the difference between hot updates and hard reloads.
    \item Discuss configurable backends, prompts, ontology sources, worker settings, tracking settings, and pipeline presets.
\end{itemize}

\section{Observability, Persistence, and Operator Control}
\begin{itemize}
    \item Describe the dashboard as a practical engineering component, not only a debugging aid.
    \item Cover live image display, scene graph view, conversation view, memory view, and configuration editing.
    \item Discuss optional persistence of the last state and image.
\end{itemize}



\section{NAOqi Application Structure}
\begin{itemize}
    \item Introduce the Pepper client service and its lifecycle.
    \item Describe how the application initializes adapters, loads configuration, and exposes robot methods such as \texttt{look}, \texttt{scan}, \texttt{ask}, and \texttt{resetConversation}.
    \item Mention the practical consequence of the client being implemented in Python 2 within the robot ecosystem.
\end{itemize}

\section{Turn Management and Interaction Logic}
\begin{itemize}
    \item Describe the \texttt{TurnManager} as the coordination layer for user-visible interaction.
    \item Explain busy-state protection, acknowledgement speech, thread-based turn execution, and fallback behavior under failure.
    \item Describe each turn type separately:
    \begin{itemize}
        \item \texttt{look}: capture one frame, request a caption, optionally trigger background detection;
        \item \texttt{scan}: move the head over predefined yaw angles, process multiple frames, optionally summarize the scan through chat;
        \item \texttt{ask}: optionally refresh visual context, then run grounded chat;
        \item \texttt{refreshAndAsk}: force a visual refresh before answering.
    \end{itemize}
\end{itemize}

\section{Visual Acquisition}
\begin{itemize}
    \item Describe the camera adapter built on ALVideoDevice.
    \item Explain frame capture, JPEG encoding, field-of-view retrieval, and timing metadata.
    \item Mention the fake-camera path used for testing and simulation.
\end{itemize}

\section{Robot Context Collection}
\begin{itemize}
    \item Explain the role of the \texttt{RobotContextCollector}.
    \item Describe how pose, people perception, face detection, social signals, and sonar measurements are collected.
    \item Cover the specific robot-side sources:
    \begin{itemize}
        \item people IDs, yaw, pitch, and distance from ALPeoplePerception;
        \item face labels and confidence from ALFaceDetection;
        \item gaze, engagement, waving, sitting, age, gender, smile, and expression cues from NAOqi social modules;
        \item left and right sonar measurements.
    \end{itemize}
\end{itemize}

\section{Metadata Construction and Session State}
\begin{itemize}
    \item Describe the metadata payload sent together with each image.
    \item Explain frame IDs, scan IDs, capture mode, camera geometry, and attached people/social information.
    \item Describe the local session store: chat ID, last caption, last detection timestamp, selected output language, and server base URL.
\end{itemize}

\section{Robot-Side Transport, Speech, and Tablet Output}
\begin{itemize}
    \item Describe the HTTP transport layer and its request types for caption, detect, chat, conversation reset, and config patching.
    \item Explain timeout and malformed-response handling.
    \item Describe speech synthesis and animated speech fallback.
    \item Describe the tablet adapter used to display the server dashboard on Pepper.
\end{itemize}

\section{Client-Side Limitations and Trade-Offs}
\begin{itemize}
    \item Discuss Python 2 compatibility, request blocking, network dependence, and limited local reasoning.
    \item Explain why the client focuses on sensing, metadata packaging, and interaction orchestration rather than heavy intelligence.
\end{itemize}

\chapter{Server-Side Integration}

\section{Conversation State Management}
\begin{itemize}
    \item Describe the conversation service and its bounded conversation history.
    \item Explain storage of original text, model text, detected language, output language, and translation status.
    \item Explain why conversation state is separate from scene memory.
\end{itemize}

\section{Prompt Composition from Scene Memory}
\begin{itemize}
    \item Describe how the chat service serializes the current world state into textual context.
    \item Explain object ordering by salience, with emphasis on socially relevant people.
    \item Describe the use of latest caption and recent captions as additional grounding context.
    \item Explain how the current implementation realizes grounded retrieval by rendering structured memory into prompts.
\end{itemize}

\section{General Scene-Aware Chat}
\begin{itemize}
    \item Describe the normal chat endpoint for open questions about the current scene.
    \item Explain language normalization, history inclusion, prompt rendering, and provider-agnostic text generation.
    \item Discuss expected question types and answer scope.
\end{itemize}

\section{Object-Focused Dialogue}
\begin{itemize}
    \item Present the object-chat mode as one of the strongest concrete contributions of the current implementation.
    \item Explain object matching by label, ranking by salience, and extraction of per-object attributes and relations.
    \item Describe the crop-caption fallback used when an object has weak structured memory but a tracked crop is available.
    \item Explain how matched object IDs are returned as traceability metadata.
\end{itemize}

\section{Direct Vision Chat}
\begin{itemize}
    \item Describe the route that bypasses scene-memory grounding and sends an image directly to a multimodal model together with a question.
    \item Explain how this mode complements the memory-grounded chat mode and when it is useful.
    \item Discuss its trade-off: richer direct visual reasoning but weaker persistent identity tracking.
\end{itemize}

\section{Multilingual Handling}
\begin{itemize}
    \item Explain the translation layer and the use of canonical model-facing language when needed.
    \item Describe the current English/Czech output modes and how translation is applied to both user input and model response.
    \item Discuss translation-related failure modes and semantic drift.
\end{itemize}

\section{Failure Handling and Practical Limitations}
\begin{itemize}
    \item Describe response fallback behavior under timeouts, malformed outputs, unavailable services, or missing visual context.
    \item Discuss remaining limits of prompt-only grounding, short conversation memory, and scene-memory incompleteness.
    \item Explicitly identify which limitations belong to the current implementation and which are reserved for future work.
\end{itemize}


\section{Role of the \texttt{research} Package}
\begin{itemize}
    \item Explain that the repository contains both production code and research tooling.
    \item Clearly separate the newer structured experiment pipeline from older exploratory scripts and notebooks.
    \item State which parts directly support the thesis evaluation.
\end{itemize}

\section{Structured Experiment Pipeline}
\begin{itemize}
    \item Present the newer \texttt{research.experiments} pipeline as the primary evaluation framework.
    \item Describe its phases in execution order:
    \begin{itemize}
        \item image description generation;
        \item vocabulary mining;
        \item draft scene graph generation;
        \item context-rot experiments;
        \item scene graph evaluation.
    \end{itemize}
    \item Explain the purpose of run metadata, stage metrics, and artifact saving.
\end{itemize}

\section{Reuse of Production Components in Experiments}
\begin{itemize}
    \item Describe the research adapters that import and reuse production detection, caption, LLM, and VLM code.
    \item Explain why this is important scientifically: the offline experiments should evaluate the same or very similar components as the deployed system.
\end{itemize}

\section{Description Generation and Vocabulary Mining}
\begin{itemize}
    \item Describe the description phase: detect objects, then generate detailed captions focused on those detections.
    \item Describe the vocabulary phase: extract candidate predicates and attributes per image, then consolidate them into a final vocabulary.
    \item Explain why vocabulary design is a key variable for scene graph quality and dialogue usefulness.
\end{itemize}

\section{Draft Scene Graph Generation}
\begin{itemize}
    \item Describe the construction of SoM-marked images for offline scene-graph prompting.
    \item Explain prompt rendering from detected objects, vocabulary, and captions.
    \item Describe the use of structured output schemas for draft graph generation.
\end{itemize}

\section{Context-Rot and Sensitivity Experiments}
\begin{itemize}
    \item Explain the idea behind context-rot: reducing available vocabulary and measuring degradation in graph output.
    \item Describe how vocabulary slices are built and evaluated.
    \item Explain how prompt sensitivity is later summarized across runs.
\end{itemize}

\section{Scene Graph Evaluation Metrics}
\begin{itemize}
    \item Define the core metrics implemented by the code:
    \begin{itemize}
        \item strict triplet precision, recall, and F1;
        \item binary relation triplet metrics;
        \item unary attribute metrics;
        \item relation-pair metrics;
        \item predicate-only metrics;
        \item overgeneration and undergeneration rates;
        \item optional normalized graph edit distance.
    \end{itemize}
    \item Explain micro-averaging and per-predicate summaries.
    \item Describe image-potency statistics and why they help characterize scene complexity.
\end{itemize}

\section{System-Level Evaluation Protocol}
\begin{itemize}
    \item Define what should be evaluated on the deployed system itself:
    \begin{itemize}
        \item latency of captioning, detection, memory update, scene graph generation, and end-to-end chat;
        \item stability of tracked identities across frames and scans;
        \item usefulness of Pepper social fusion for person-centered dialogue;
        \item qualitative quality of grounded robot responses in interactive scenarios.
    \end{itemize}
    \item Explain how pipeline stage timings exposed by the server can be used for runtime profiling.
\end{itemize}

\section{Legacy Exploratory and Training Pipelines}
\begin{itemize}
    \item Describe the older exploratory workflows, but keep them clearly secondary to the main evaluation story.
    \item Mention:
    \begin{itemize}
        \item automated SoM scene-graph dataset generation with Grounding DINO and an LLM labeler;
        \item knowledge distillation for detector training data;
        \item detector training for RT-DETR or YOLO;
        \item LoRA-based VLM fine-tuning with Unsloth.
    \end{itemize}
    \item Explain whether these are reported as completed experiments, supporting tooling, or future-ready extensions.
\end{itemize}

\chapter{Robot-Side Integration}

\section{FastAPI Application and Service Composition}
\begin{itemize}
    \item Describe the server startup process, application state, and service initialization.
    \item Explain how configuration loading leads to either an in-process perception pipeline or worker-backed runtime.
    \item Present the main route groups: detection, captioning, chat, vision chat, memory, config, worker, and dashboard.
\end{itemize}

\section{Detection Subsystem}
\begin{itemize}
    \item Describe the detection service abstraction and the supported backends: YOLO, RT-DETR, RF-DETR, and OWL-V2.
    \item Explain ontology handling for open-vocabulary detection.
    \item Discuss thresholding, device selection, and detector interchangeability.
\end{itemize}

\section{Persistent Tracking and Identity Maintenance}
\begin{itemize}
    \item Describe the role of \texttt{SceneMemory} in maintaining object continuity across frames.
    \item Explain appearance embedding extraction, crop preparation, and use of a visual encoder for re-identification.
    \item Explain the Hungarian matching procedure and the weighted combination of appearance and geometry.
    \item Cover unmatched-track aging, dormant-track removal, and object hit statistics.
\end{itemize}

\section{Pepper-Specific Fusion with Robot-Native People Perception}
\begin{itemize}
    \item Explain why server-side visual detections alone are insufficient for socially grounded interaction.
    \item Describe geometric projection from robot person angles into image coordinates.
    \item Explain the matching between Pepper people and server person detections using angular similarity, previous bindings, tracker stability, and distance-size consistency.
    \item Describe the generation of synthetic person detections when Pepper reports a person but no visual bounding box matches.
    \item Explain how social attributes such as waving, sitting, looking at the robot, engagement zone, age, gender, smile, and expressions are merged into object state.
\end{itemize}

\section{Set-of-Mark Rendering}
\begin{itemize}
    \item Describe why Set-of-Mark visualization is used for VLM grounding.
    \item Explain bounding boxes, labels, optional masks, optional polygons, and style scaling.
    \item Mention the available mask backends and the role of the overlay in the scene graph prompt pipeline.
\end{itemize}

\section{Scene Graph Generation}
\begin{itemize}
    \item Present the scene graph service as a strategy layer with four modes: rules, VLM, RelTR, and hybrid.
    \item Explain the rule-based backend:
    \begin{itemize}
        \item configurable predicates;
        \item spatial, directional, overlap, containment, and label-pair rules;
        \item automatic color-attribute extraction.
    \end{itemize}
    \item Explain the VLM backend:
    \begin{itemize}
        \item prompt construction from predicates and captions;
        \item structured output parsing and repair;
        \item filtering of hallucinated relations against current detections.
    \end{itemize}
    \item Explain the optional RelTR backend and why it is treated as a complementary relation source.
    \item Explain the hybrid mode as a union of complementary evidence sources.
\end{itemize}

\section{Scene Memory as a Dynamic World Model}
\begin{itemize}
    \item Describe the internal state representation: tracked objects, relationships, captions, Pepper-person bindings, and timestamps.
    \item Explain object fields that matter for dialogue: label, attributes, status, hits, frame IDs, scan IDs, estimated bearing, Pepper person ID, and robot distance.
    \item Explain relationship accumulation, unary-attribute handling, and caption insertion.
    \item Cover memory pruning, TTLs, and configurable object, relation, and caption limits.
\end{itemize}

\section{Captioning and Auxiliary Vision Services}
\begin{itemize}
    \item Describe the captioning subsystem and its role in fast single-frame scene summarization.
    \item Explain that the caption endpoint can launch detection asynchronously so that later dialogue is grounded in updated memory.
    \item Describe the separate \texttt{vision\_chat} endpoint for direct image-plus-question multimodal interaction.
\end{itemize}

\section{Worker Runtime and Runtime Switching}
\begin{itemize}
    \item Describe the worker manager, internal worker API, and lazy startup behavior.
    \item Explain why the worker mode exists: isolate heavy inference, manage warmup, and keep the web process lighter.
    \item Cover health checks, restart logic, circuit breaker behavior, and hot runtime updates.
\end{itemize}

\section{Dashboard and Runtime Administration}
\begin{itemize}
    \item Describe the dashboard pages and their practical roles.
    \item Explain live WebSocket updates for detections, captions, memory, and chat messages.
    \item Cover operator-editable settings: detector backend, thresholds, scene graph mode, VLM prompts, worker controls, memory limits, chat/caption providers, and pipeline presets.
\end{itemize}

\chapter{Experimental Methodology}

\section{Perception and Scene-Memory Results}
\begin{itemize}
    \item Report detector behavior, tracking behavior, and Pepper-fusion behavior.
    \item Discuss what types of errors dominate: missed detections, identity fragmentation, false relations, or noisy social attributes.
    \item Show examples where scene memory improves follow-up interaction.
\end{itemize}

\section{Scene Graph Quality Results}
\begin{itemize}
    \item Report the graph metrics defined in the previous chapter.
    \item Compare different prompting setups, model providers, vocabularies, or graph-generation modes if these experiments are executed.
    \item Discuss the balance between richer semantics and hallucination risk.
\end{itemize}

\section{Dialogue and Interaction Results}
\begin{itemize}
    \item Present representative interaction transcripts for \texttt{look}, \texttt{scan}, general chat, object chat, and vision chat.
    \item Analyze whether answers remain grounded in the scene and whether the robot maintains object continuity across follow-up questions.
    \item Discuss multilingual behavior and user-facing latency.
\end{itemize}

\section{Ablations and Sensitivity Findings}
\begin{itemize}
    \item Interpret vocabulary-size sensitivity, prompt sensitivity, and any comparisons between rule-based and multimodal graph generation.
    \item Discuss which system parts most strongly affect downstream dialogue quality.
\end{itemize}

\section{Architectural Discussion}
\begin{itemize}
    \item Reflect on the main architectural choices:
    \begin{itemize}
        \item remote inference versus onboard intelligence;
        \item in-process versus worker runtime;
        \item explicit scene memory versus implicit conversational context;
        \item hybrid graph generation versus a single-model solution.
    \end{itemize}
    \item Explain what worked well and what remains engineering debt.
\end{itemize}

\section{Limitations}
\begin{itemize}
    \item Current reliance on external compute and network connectivity.
    \item Limited long-term memory and lack of a richer knowledge layer beyond current scene state and short conversation history.
    \item Dependence on prompt quality and provider reliability.
    \item Incomplete evaluation of some exploratory training paths if they remain unfinished.
    \item Potential mismatch between offline graph metrics and perceived dialogue quality in real HRI.
\end{itemize}

\chapter{Results, Discussion, and Conclusion}

\section{Summary of the Thesis}
\begin{itemize}
    \item Recap the problem, the implemented system, and the main findings.
    \item Summarize how the thesis connects Pepper sensing, server-side perception, dynamic scene memory, and grounded dialogue.
\end{itemize}

\section{Main Contributions Revisited}
\begin{itemize}
    \item Restate the final contributions in concise, evidence-backed form.
    \item Make sure this section reflects what was actually implemented and evaluated, not only initial intentions.
\end{itemize}

\section{Future Work}
\begin{itemize}
    \item Stronger long-term memory and user modeling.
    \item Better open-vocabulary perception and relation modeling.
    \item More principled multimodal grounding and uncertainty estimation.
    \item Additional user studies and broader HRI evaluation.
    \item Tighter integration with other Pepper behaviors and action-capable planners.
\end{itemize}

%% Temporary drafting state: print the full bibliography while shaping the related-work chapter.
\nocite{*}
\printbibliography[heading=bibintoc]

\end{document}